\originalchapter{语言模型的数据与可信评估}\label{ch:learning}
\originalsection{学到的是哪一种分布}
优化解决了如何降低一个确定的目标函数. 但语言模型在训练语料上降低损失, 并不等于能够回答新的问题. 一份被记住的试卷和一份从未见过的试卷可能有完全不同的表现. 本章把“训练时优化什么”与“用什么证据判断学到了什么”连起来. 这些区别在预训练、指令微调和后训练中都存在, 只是样本单位、损失和评价任务有所不同.

\begin{definition}
设 $Z$ 是一个按约定抽取的评价单位, 其分布为 $P$. 对模型 $f$ 和损失 $\ell(f,Z)$, 总体风险为 $R_P(f)=\E_{Z\sim P}\ell(f,Z)$. 若 $Z_1,\ldots,Z_n$ 为实际观测, 经验风险为
\[
\widehat R_n(f)=\frac1n\sum_{i=1}^n\ell(f,Z_i).
\]
这里的 $f$ 是包含参数和使用规则的整个预测系统. 提示词、检索库、采样温度和工具使用规则若影响输出, 也属于评价时须固定的设计.
\end{definition}
总体风险是生成过程的平均, 往往不能直接算出; 经验风险是有限数据的平均, 可以计算. 当算法根据数据选择模型 $\widehat f$, 这个模型也是一个随机对象. 因而对事先固定的 $f$ 成立的统计结论, 不能直接把 $f$ 换成在同一数据上挑出的 $\widehat f$. 后面会精确说明何时可以这样替换.

语言模型的一个文档是序列 $x_1,\ldots,x_T$. 自回归模型按链式分解给出
\[
p_\theta(x_1,\ldots,x_T)=\prod_{t=1}^T p_\theta(x_t\mid x_{<t}),\qquad
-\log p_\theta(x_1,\ldots,x_T)=\sum_{t=1}^T-\log p_\theta(x_t\mid x_{<t}).
\]
第一个 token 的上下文可以包含固定的起始符. 等式来自条件概率的乘法规则, 不是神经网络专有的假设. 模型决定的是每个条件概率怎样由上下文计算. 预训练通常优化大量实际文档中的这些负对数概率之和或平均.

平均的分母须明确. 对每篇文档先算平均 token 损失, 再对文档平均, 是让每篇文档具有相同权重; 把所有文档的 token 损失相加再除以 token 总数, 是让长文档获得更多权重. 二者对应不同的目标总体. 若研究“随机用户提交一篇文档的表现”, 文档平均可能更贴近问题; 若研究语料压缩或随机 token 的预测, token 平均通常更自然. 不能在比较模型时暗中更换分母.

\begin{example}
两篇文档分别含 2 个与 8 个评价 token, 每个 token 的平均负对数似然分别为 1 与 3. 求文档平均和 token 平均.
\end{example}
\begin{solution}
文档平均为 $(1+3)/2=2$; token 平均为 $(2\cdot1+8\cdot3)/10=2.6$. 第一种抽样先等概率选文档, 第二种抽样在 10 个 token 中等概率选一个. 两个计算都正确, 但回答的问题不同. 若长文档恰好更难, 只改变权重就能改变报告值, 无须模型行为改变.
\end{solution}

\originalsection{数据划分、污染与开发过程}
训练集用于更新参数, 开发集或验证集用于选择设计, 测试集用于设计确定后的最终评价. 预训练数据、后训练数据和评价数据还须区分: 某个基准没有进入指令微调集, 不能据此断定它没有进入预训练语料. 题面、答案、解析、改写和同源版本都可能传递评价信息. 这种污染可以造成较高分数, 但不能仅凭分数高就断定存在污染.

先确定部署问题, 才能决定怎样划分. 如果目标是预测已有作者的新文章, 按文档留出可能合理; 如果目标是新作者, 须按作者分组; 如果目标是将来发生的知识和事件, 须考虑时间切分. 按时间切分仍不能自动排除转载或提前公开的答案. 划分只是评估设计的一部分, 必须配合来源与去重检查.

语言模型常把文档切为长度固定的窗口. 若先切窗口再随机划分, 同一文档的相邻窗口可能跨过训练与测试边界. 训练窗口的后半段还可能与测试窗口的前半段完全相同. 此时测试不再模拟新文档. 更合理的顺序是先在文档、用户或其他目标单位上划分, 再在各集合内部切窗口. 精确去重可比较内容哈希; 近重复需要另定比较规则, 并记录阈值、误删风险与无法识别的改写.

开发集也会被逐渐学会. 假设每次试验都看同一批 32 道题, 再改变提示词、样例排列、停止规则或模型版本. 最终设计不只是由训练参数决定, 还受这 32 道题影响. 即使没有对开发集反向传播, 这种人为选择也构成信息使用. 有效做法是记录尝试次数与选择准则, 在选择结束后使用独立留出集, 或在外层留出数据而只在内层选择设计. 数据少时可以交叉验证设计的稳定性, 但不能把反复查看外层结果仍称为独立测试.

\begin{exercise}
从 20 篇文档各切出 100 个重叠窗口, 随机按 8:2 分到训练和测试. 另有一套考试题, 其标准答案未入训练集, 但详细解析入了训练集. 分别说明两个设计的风险, 并提出修正.
\end{exercise}
\begin{solution}
重叠窗口会分享文本和文档背景, 所以 2000 个窗口不等于 2000 个独立的新文档. 应先按文档或部署所需的作者、时间单位划分, 再切窗口. 第二种情形虽然未见标准答案文件, 解析仍可能包含答案及解法信息, 因而不能作为未见题泛化的证据. 应把同源题面、答案和解析放在同一侧, 检查近重复与改写, 再选择来源和时间清楚的留出题. 如果无法核实预训练语料, 就应报告污染状态未知, 不把它写成已排除.
\end{solution}

\originalsection{预训练数据的受控比较}
评价不仅用于比较模型, 也用于研究数据处理是否改善学习. 网页中的菜单、重复模板和大量镜像文本可能占据训练预算; 过滤与去重改变了同样预算里模型实际看到的内容. 但过滤也可能删去罕见语言、专业文本或与部署有关的格式, 因而“保留得更少”并不自动等于“数据质量更高”.

FineWeb 的数据消融提供了具体的阅读对象\cite{fineweb}. 它在同一比较中固定模型架构、参数规模、训练 token 数和评价设置, 主要改变文本抽取、过滤或去重步骤; 使用两组随机数据子集和初始化来降低单次实验偶然性. 相关消融采用约 1.71B 参数的模型, 过滤研究常使用约 28B 训练 token, 一些去重与累计比较使用 350B token. 读者应先核对是哪一种预算, 再阅读分数, 不能把不同预算的曲线直接当作纯数据效应.

其中一个有教学价值的结果是: 逐快照去重与跨全部快照的全局去重并不等价, 后者在所研究设置中没有产生原先预期的明显改善. 去重也改变了内容混合比例, 因而“删除重复”不能单独解释所有分数变化. 这里的实测支持是在给定模型、预算和任务组中, 某些处理方案比对照更好; 对哪些内容促成了何种内部机制, 仍需额外实验. 这既不证明所有数据集都应照搬同一过滤规则, 也不证明处理步骤改善了全部部署任务.

读一项数据论文时, 可以提出可回答的问题: 比较中训练 token 总数是否相同? 删除后是否重新采样而提高某类内容权重? 是否同时更换了分词或模型? 评价任务是否也参与筛选数据规则? 这些问题把“更好的数据”变成可检验的具体干预. 若要在本书 CPU 规模上练习, 可固定一个小模型与更新预算, 比较含重复模板和去重后的合成序列, 预先声明用哪些未见组合计分. 这只是方法迁移的设计, 本章没有把它冒充已经复现的 FineWeb 实验.

\originalsection{有限样本能够保证什么}
先看一个可完全证明的有限模型结论. 它适用于落在 $[0,1]$ 的损失, 如答题错误指示量, 或事先定义的文档级有界评分. 原始负对数似然没有这样的统一上界: 当正确 token 概率趋于零, $-\log p$ 可以任意大. 不能把下面的结论原样用于未截断的语言模型交叉熵. 若预先把损失截断为 $\min(\ell,L)/L$, 界约束的是截断后的风险; 极低概率事件的尾部损失仍需另行研究.

\begin{lemma}[有界随机变量的指数估计]
若 $0\le Z\le1$, 对任意实数 $t$, 有 $\E e^{t(Z-\E Z)}\le e^{t^2/8}$.
\end{lemma}
\begin{proof}
记 $F(t)=\log\E e^{tZ}$. 因 $Z$ 有界, 在任何有限 $t$ 区间上, 指数及其前两阶导数都由一个有限常数控制, 可在期望内求导. 计算得
\[
F'(t)=\frac{\E[Ze^{tZ}]}{\E[e^{tZ}]},\qquad
F''(t)=\frac{\E[Z^2e^{tZ}]}{\E[e^{tZ}]}-
\left(\frac{\E[Ze^{tZ}]}{\E[e^{tZ}]}\right)^2.
\]
把原分布按正权重 $e^{tZ}/\E e^{tZ}$ 重新加权, 这仍是一个概率分布, 且仍满足 $0\le Z\le1$. 上式是该分布中的方差. 对任意取值于 $[0,1]$ 的变量, 若均值为 $\mu$, 由 $Z^2\le Z$ 得方差至多 $\mu-\mu^2\le1/4$. 因而 $F''(t)\le1/4$.

令 $G(t)=F(t)-t\E Z$, 则 $G(0)=G'(0)=0$. 二次积分余项写为
$G(t)=t^2\int_0^1(1-s)F''(st)\,ds$, 对正负 $t$ 都成立. 所以 $G(t)\le t^2/8$. 对两侧取指数即得结论.
\end{proof}

\begin{theorem}[固定模型的 Hoeffding 界]
若 $Z_i$ 独立同分布, 固定模型的损失 $L_i=\ell(f,Z_i)$ 落在 $[0,1]$, 则对任意 $\epsilon>0$,
\[
\Pr\left(\left|\widehat R_n(f)-R_P(f)\right|\ge\epsilon\right)
\le 2e^{-2n\epsilon^2}.
\]
\end{theorem}
\begin{proof}
设 $R=\E L_i$. 对任意 $t>0$, 事件 $\sum_i(L_i-R)\ge n\epsilon$ 蕴含 $e^{t\sum_i(L_i-R)}\ge e^{tn\epsilon}$. 非负变量的期望至少为这个阈值乘事件概率, 从而
\begin{align*}
\Pr\left(\sum_i(L_i-R)\ge n\epsilon\right)
&\le e^{-tn\epsilon}\E e^{t\sum_i(L_i-R)}\\
&=e^{-tn\epsilon}\prod_i\E e^{t(L_i-R)}
\le e^{-tn\epsilon+nt^2/8}.
\end{align*}
独立性用在把期望拆成乘积这一步. 右边指数是关于 $t$ 的二次式, 取 $t=4\epsilon$ 得 $e^{-2n\epsilon^2}$. 对负偏差应用引理中的 $-t$, 得同样上界. 正负两事件的并集概率至多为概率之和, 所以乘以 2.
\end{proof}

\begin{proposition}[有限候选的同时保证]
设 $f_1,\ldots,f_M$ 是不依赖本次评价数据的固定候选, 均使用 $[0,1]$ 损失. 对独立同分布的 $n$ 个评价单位及 $0<\delta<1$, 至少以概率 $1-\delta$, 所有候选同时满足
\[
|\widehat R_n(f_j)-R_P(f_j)|\le
\epsilon_n,\qquad \epsilon_n=\sqrt{\frac{\log(2M/\delta)}{2n}}.
\]
在此事件上, 令 $\widehat f$ 最小化候选的经验风险, $f_*$ 最小化候选的总体风险, 则 $R_P(\widehat f)\le R_P(f_*)+2\epsilon_n$.
\end{proposition}
\begin{proof}
每个候选失败的概率至多为 $2e^{-2n\epsilon_n^2}$. 任意一个失败的概率至多为它们之和 $2Me^{-2n\epsilon_n^2}=\delta$; 无须假定不同模型的损失相互独立. 在所有界成立时,
\[
R_P(\widehat f)\le\widehat R_n(\widehat f)+\epsilon_n
\le\widehat R_n(f_*)+\epsilon_n\le R_P(f_*)+2\epsilon_n.
\]
中间一步来自经验最小化, 两端来自同时偏差界. 因为已经同时控制所有候选, 数据选出的下标也在控制范围内.
\end{proof}
若 $M=1000,n=10000,\delta=0.05$, 得 $\epsilon_n\approx0.0230$. 这是充分保证, 通常很保守, 不是实际误差必然达到 2.3 个百分点. 若一次评价单位是一篇文档, $n$ 就是文档数, 不能擅自改成相关 token 的总数. 候选网络由独立训练集产生时, 可以先条件于训练集, 再对独立开发数据应用此界. 若新的候选根据此前的开发反馈逐步产生, 便不满足“本次评价前已固定”的条件, 不能简单以实际尝试次数代入宣称获得同样保证.

最终模型若只由训练和开发数据选出, 则条件于这两部分数据后, 它对独立测试集是固定模型, 可应用固定模型界. 这解释了保持测试集独立的作用. 它没有证明测试分布与部署分布相同, 也没有处理评分器本身的偏差.

\originalsection{正则化与选择偏差的实测}
语言模型在大量参数和有限样本之间需要某种选择约束. 例如目标
$\widehat R_n(\theta)+\lambda\norm{\theta}^2/2$ 对参数大小施加惩罚, 其梯度比原损失增加 $\lambda\theta$. 早停限制更新次数, 数据增强规定哪些变化应保留答案, 参数高效微调限制可更新的方向. 它们改变训练算法或可选择的函数, 但不自动证明模型学到了因果特征. 惩罚系数、停止步数和提示方案同样需要开发数据选择. 用测试分数选择它们仍会造成选择偏差.

新增 \texttt{v03\_networks.py} 的 \texttt{selection} 实验只隔离评价中的选择效应. 它生成已固定候选的独立有界损失, 每个候选的真实错误风险都是 $0.5$. 每次只用 32 个验证单位挑选经验错误率最小的候选, 再独立抽取 2000 个测试损失来评价选中者. 400 次重复中, 候选数为 $1,8,64,512$ 时, 选中者验证错误率的平均依次为 $0.49969,0.37117,0.29766,0.23750$; 独立测试平均依次为 $0.50034,0.49986,0.49925,0.50013$.

验证分数随着搜索规模提升, 真实能力完全没有提升. 原因不是优化器有问题: 从多个带噪声的平均里挑最小值, 同时也挑中了偏低的采样噪声. 这是合成损失表, 没有训练神经网络; 其用途是显示机制, 不能把 26.3 个百分点的乐观差距当成现实基准污染量. 完整指标、随机种子 144、环境和源码哈希保存在 \texttt{results/v03-networks/}.

\originalsection{概率、校准与任务指标}
语言模型的输出是概率分布, 但通常还要转成回答、程序或动作. Token 负对数似然度量模型给真实后继分配多少概率; 答题正确率度量按约定生成与评分后的答案. 同一个事实可以有多种正确表达, 所以真实文本概率不等于答案正确率. 相反, 在有限题目上猜中答案也不能证明概率分布拟合良好. 工具系统还要区分提案、执行和最终完成, 检索系统要区分检索命中和基于证据的回答.

若使用自然对数, 困惑度定义为平均 token 负对数似然的指数. 在所有位置都给正确 token 相同概率 $p$ 的特殊情形, 困惑度为 $1/p$. 一般情形是这些倒数的几何平均. 比较困惑度须使用相同的分词、评价文本和计分位置; 词表不同会改变一个 token 所含的信息量. 不能只比较数值就判断不同分词模型谁更好.

\begin{definition}
对二分类概率预测 $q(X)\in[0,1]$, 校准指 $\E[Y\mid q(X)]=q(X)$. 意思是同样预测值的一组样本中, 正类的实际频率与该预测值相符. 实际样本有限, 通常把概率分箱后比较各箱的平均预测和正类频率, 而不是声称逐个实数概率都已验证.
\end{definition}
准确率只看阈值两侧, 无法识别过度自信. 一个正确率为 80\% 的系统若对每个答案都报告 99\% 的信心, 就可能严重失准. 语言模型还须明确“信心”指哪一个事件: 单 token 的概率、完整字符串的概率、某答案集合的总概率, 或另一个评分器给出的成功概率, 这些不是同一个量.

平方概率误差称为 Brier 损失. 若真实条件概率为 $p=\Pr(Y=1\mid X=x)$, 则
\[
\E[(q-Y)^2\mid X=x]=p(q-1)^2+(1-p)q^2
=(q-p)^2+p(1-p).
\]
因而在这个固定 $x$ 上, 最小值在 $q=p$ 处取得. 二分类对数损失的条件期望为 $-p\log q-(1-p)\log(1-q)$, 对 $0<p<1$ 求导得 $(q-p)/(q(1-q))$, 同样在 $q=p$ 处取得最小值. 这说明两种损失都鼓励报告真实概率. 但有限数据、训练误差和分布偏移都可能使实际模型不校准; 常数预测等于总体正类率可以校准, 却几乎不区分不同输入.

一个简单的后处理是温度缩放: 对分类 logits $z$, 使用 $\softmax(z/\tau)$, 其中 $\tau>0$ 在独立校准集上选择\cite{calibration}. 正温度不改变最大 logit 的位置, 只改变概率集中程度. 它因此可以改善概率指标而保持按最大概率选出的类别不变, 也无法修复分类排序错误. 温度过大趋向均匀, 过小趋向把概率集中到最大值. 每次修改温度仍是模型选择, 不应依据最终测试集调整.

新增 \texttt{calibration} 实验令 $X\sim N(0,1.5^2)$, $Y\mid X\sim\operatorname{Bernoulli}(\sigma(X))$, 并故意给模型 logits $3X$. 它使用 800 个验证样本, 在 $0.5$ 到 $5$ 的 91 个正温度中以验证负对数似然选择温度, 然后只评价一次 5000 个独立测试样本. 种子 144 得温度 $2.9$, 接近生成规则中的 3. 这不是训练得到的现实模型, 是一个可控制的过度自信案例. 测试准确率前后均为 $0.7362$, 测试对数损失从 $0.72520$ 降至 $0.52642$, Brier 损失从 $0.19905$ 降至 $0.17577$. 分箱误差也记录下来, 但它依赖箱数与采样, 并不等于校准定义本身. 脚本计算的是正类概率的分箱误差, 不是多分类最大类别信心的 ECE.

\originalsection{分布外表现与因果限制}
设训练和测试来自 $P$, 部署来自 $Q$. 上面的概率界只控制 $R_P$, 不控制 $R_Q$. 例如训练文档中的某个来源标记总与标签共现, 模型可以依靠标记取得极低经验错误率. 当部署中标记与标签反相关, 原有规律失效. 对基础模型, 新领域、语言、上下文长度、时间和工具环境都可能改变分布; “分布外”必须说明相对哪个训练或评价分布而言.

旧 \texttt{spurious} 实验仍提供一个明确算例: $s=2y-1$, 稳健特征 $x_1=s+0.9\varepsilon_1$, 快捷特征在训练和同分布测试为 $x_2=s+0.1\varepsilon_2$, 偏移测试变为 $x_2=-s+0.1\varepsilon_2$. 独立标准正态噪声使第二个特征在训练中更容易使用. 保留 $x_1$ 的对照是作者知道生成规则后设计的干预, 不是能自动从未知语言数据中找出因果特征的算法. 语言模型的类似问题可以通过改变来源标记、模板、实体组合或检索证据来构造, 但每个改动都要核对是否改变了正确答案.

相关性只说明在观察到的分布中共同变化. 要研究“模型是否依据证据回答”, 可以保持问题不变, 有控制地替换支持证据并预测答案应怎样改变; 要研究组合泛化, 可以把某些语法组合留出, 同时让相关词和基本规则在训练中出现. 这些设计比只报告一个总分更接近机制问题. 但替换可能引入不自然文本, 留出可能同时改变难度, 评分器也可能出错. 因而需要对照、错误分析和明确样本量, 不能把一次干预等同于完整因果解释.

\begin{exercise}
一个模型在 100 道题上答对 90 道. 研究者测试 200 个提示词并报告最好提示的 90\% 正确率, 再用这 100 道题选择温度. 能否将其当作独立测试正确率与校准结果? 如何修正?
\end{exercise}
\begin{solution}
不能. 题目已经同时参与提示选择和概率后处理选择, 实际是开发数据. 可以固定已选提示和温度, 在新的同分布题目上评价, 同时另设目标明确的分布外题组. 若 200 个提示在评价前全部固定, 有限模型界可同时控制这些提示在有界评分上的偏差, 但它仍依赖独立评价单位, 且不给校准定义的完整保证. 若提示是根据这 100 题的失败逐步改写, 简单代入 $M=200$ 不足以恢复已被破坏的条件.
\end{solution}
\begin{exercise}
二分类模型在所有输入上都预测 $q=0.2$, 总体正类率恰为 $0.2$. 它是否校准? 这说明了什么? 若部署正类率变为 $0.6$, 又如何?
\end{exercise}
\begin{solution}
原分布下 $\E[Y\mid q=0.2]=0.2$, 因而校准. 但它没有根据输入区分不同事件, 所以校准不等于有用的表征或高区分能力. 部署分布下同一预测组的正类率为 $0.6$, 于是不校准. 在一个分布上通过校准检查, 不保证在另一分布上仍校准.
\end{solution}
\begin{exercise}
对 $M=20$ 个事先固定的候选, 需要以至少 $0.95$ 的概率让有界风险偏差均不超过 $0.05$. 使用本章结论给出充分样本数, 并说明它没有保证什么.
\end{exercise}
\begin{solution}
要求 $2M e^{-2n(0.05)^2}\le0.05$, 得
$n\ge\log(800)/0.005\approx1336.923$, 所以取 $n=1337$ 个独立评价单位足够. 它不保证所有候选的风险都低, 不保证选中模型是全体可能模型中最优, 不保证未截断对数损失, 也不保证部署分布不同后的表现. 条件及结论范围是这项计算的一部分.
\end{solution}
\label{ch:end:05}
